\section*{Building from Source}
\addcontentsline{toc}{section}{Building from Source}

\subsection*{Build System Architecture}

Symphony's build system reflects our research priorities and architectural innovations. We employ a multi-stage build process that optimizes for both development velocity and production performance, enabling rapid iteration during research phases while maintaining deployment reliability.

The build architecture supports our dual ensemble approach, with separate compilation paths for Rust infrastructure components and Python orchestration systems. This separation enables independent optimization and facilitates comparative performance analysis across different architectural configurations.

\begin{table}[ht]
\centering
\begin{tabular}{@{}llll@{}}
\toprule
\textbf{Build Target} & \textbf{Technology} & \textbf{Output} & \textbf{Purpose} \\
\midrule
Core Infrastructure & Rust + Cargo & Native binaries & High-performance execution \\
Conductor System & Python + PyO3 & Extension modules & AI orchestration \\
Frontend Interface & React + Vite & Web assets & User interface \\
Extension Runtime & Multi-language & Plugin binaries & Extensibility support \\
\bottomrule
\end{tabular}
\caption{Build System Components and Targets}
\end{table}

\subsection*{Development Build Configuration}

Development builds prioritize debugging capabilities and rapid iteration over performance optimization. Our development configuration includes logging, performance profiling, and research instrumentation that enables detailed analysis of system behavior.

\begin{alertbox}
Development builds include research instrumentation that may impact performance. These builds are optimized for analysis and debugging rather than production deployment.
\end{alertbox}

Key development build features include:
\begin{compactlist}
    \item Complete debug symbol generation for all components
    \item Performance profiling integration with research analytics
    \item Hot-reload capabilities for rapid development iteration
    \item Extended logging and tracing for behavioral analysis
    \item Memory safety checks and resource usage monitoring
\end{compactlist}

\subsection*{Production Build Optimization}

Production builds implement aggressive optimization strategies that reflect our research findings on performance-critical pathways. The optimization process applies insights from our empirical analysis to achieve maximum system performance.

Our production build process includes:
\begin{expandedlist}
    \item \textbf{Link-Time Optimization}: Cross-module optimization for Rust components
    \item \textbf{Profile-Guided Optimization}: Optimization based on empirical usage patterns
    \item \textbf{Dead Code Elimination}: Removal of unused functionality to minimize binary size
    \item \textbf{Dependency Optimization}: Selective inclusion of required dependencies only
\end{expandedlist}

\subsection*{Cross-Platform Compilation}

Symphony's cross-platform build system supports our research objective of universal AI-first development environment deployment. We maintain build configurations for multiple target platforms while preserving performance characteristics across different architectures.

\begin{table}[ht]
\centering
\begin{tabular}{@{}lll@{}}
\toprule
\textbf{Platform} & \textbf{Architecture} & \textbf{Optimization Focus} \\
\midrule
Linux & x86\_64, ARM64 & Server deployment, research environments \\
macOS & x86\_64, ARM64 & Developer workstations, performance analysis \\
Windows & x86\_64 & Enterprise deployment, compatibility \\
\bottomrule
\end{tabular}
\caption{Supported Platforms and Optimization Strategies}
\end{table}

Cross-platform compilation enables comparative performance analysis across different operating systems and hardware architectures, supporting our research into platform-specific optimization strategies.
